% Part: many-valued-logic
% Chapter: syntax-and-semantics
% Section: formulas

\documentclass[../../../include/open-logic-section]{subfiles}

\begin{document}

\olfileid{mvl}{syn}{fml}

\olsection{\usetoken{P}{formula}}

\begin{defn}[Formula]
\ollabel{defn:formulas}
Himpunan~$\Frm[L]$ saka \emph{!!{formula}s} ing basa
proposisional~$\Lang L$ ditetepake kanthi induktif kaya mangkene:
\begin{enumerate}
\item Saben !!{propositional variable}~$\Obj p_i$ iku
  !!{formula} atomik.
\item Saben konektif kanthi $0$ panggonan (tetapan proposisional)
  ing~$\Lang L$ iku !!{formula} atomik.
\item Yen $\star$ konektif kanthi $n$ panggonan ing~$\Lang L$,
  lan $!A_1$, \dots, $!A_n$ iku !!{formula}s, mula
  $\star(!A_1, \dots, !A_n)$ iku !!a{formula}.
\tagitem{limitClause}{Ora ana liyane kang !!a{formula}.}{}
\end{enumerate}
Yen $\star$ nduweni $1$ panggonan, $\star(!A_1)$ asring
ditulis prasaja minangka $\star !A_1$. Yen $\star$ nduweni
$2$ panggonan, $\star(!A_1,!A_2)$ asring ditulis
minangka $(!A_1 \star !A_2)$.
\end{defn}

Kaya lumrahe, kurung paling njaba asring ora kita tulis
tanpa katrangan maneh.

\begin{ex}
  Ing basa baku~$\Lang{L_0}$, $\Obj p_1 \lif (\Obj p_1
  \land \lnot \Obj p_2)$ iku formula. Ing basa logika produk,
  formula iku ditulis dadi $\Obj p_1 \lif (\Obj p_1 \odot
  \lnot \Obj p_2)$. Yen konektif kanthi $1$ panggonan, yaiku
  $\triangle$, ditambahake ing basa, kita uga bisa duwe formula kaya
  $\triangle (\Obj p_1
  \land \Obj p_2) \lif (\triangle \Obj p_1 \land \triangle \Obj p_2)$.
\end{ex}

\end{document}
